\section{The scalar component detected by the doubling integral}

% BEGIN D87C1
\paragraph{D87C1. A cuspidal scalar form on the original principal level.}
Fix a member in the range of F86D2 with its linked realization $X$.
Use the upper definite Hermitian space $V$ of rank $n=r+t$ and the
lower space $W$ of rank $r$. Write $j$ for the model index and
$u_0$ for its tuple index; $t$ always denotes the length of the
distinguished string. The realization includes the original Weil
representation and level models ($X.\mathrm{given},X.\mathrm{full}$).
For an upper vector $y$ let $Y=\operatorname{ext}_V(y)$, with values
in its archimedean coefficient representation $(C,\rho)$.
It is a continuous function on $[U(V)]=U(V)(F)\backslash U(V)(\mathbb A_F)$.
Indeed the extension of the finite form is $\rho(h)^{-1}f_y(h_f)$;
the automorphy of the finite form cancels the change under a rational
left translation. If $b\in K=U(V)(F_\infty)$, then
$Y(hb)=\rho(b)^{-1}Y(h)$. The map $y\mapsto Y$ is injective:
the finite coordinates range over the full image of the finite
projection and $\rho(h)$ is invertible. Vanishing of $Y$ would
force $f_y=0$, hence $y=0$ by the original level identification.

Put $Z=P_g\cup Q(u_0)$, the union of the designated type places
and the twist places in $E$. Write $S$ for their restrictions to
$F$. In the given finite frame a principal-level element is integral
with integral inverse and reduces to the identity at every place of
$Z$. It belongs to the defining flag level. All determinant characters
and the tuple of upper-left blocks in its level action evaluate to
their value at the identity. Consequently every component $u=\ell Y$
is fixed by this principal level. At infinity the translates of $u$
belong to the space of functions $(\ell'Y)_{\ell'\in C^*}$.
This gives compact-type finiteness and smoothness there, since
$C$ is a finite-dimensional continuous representation of a compact
Lie group. At finite places the principal invariance gives smoothness.
The definite group has compact adelic quotient: its rational center
is the anisotropic norm-one torus and the Hermitian form has no
nonzero rational isotropic vector. There are no proper rational
parabolic constant terms to test. Thus $u$ is a smooth cuspidal
form of the stated level and is square integrable for the specified
quotient measure. These conclusions hold for every $\ell$, before
any choice of Fock test.
% END D87C1

% BEGIN D87C2
\paragraph{D87C2. Two conjugate matrix coefficients at the real places.}
Use inner products linear in their second argument. Write
$u_z(h)=\langle z,Y(h)\rangle$ for $z\in C$. For any $b\in K$
one has $u_z(hb)=u_{\rho(b)z}(h)$. Form the invariant Hermitian
pairing on $C$ by integrating $u_z\overline{u_{z'}}$ over the
quotient. It is conjugate-linear in $z$ and linear in $z'$.
Right invariance of quotient measure makes it $K$-invariant.
Irreducibility of the member's coefficient representation on $K$
and Schur's lemma make it $c_y\langle z,z'\rangle$. If $y\ne0$
there are $h,z$ with $u_z(h)\ne0$ by D87C1. Continuity and
full support of quotient Haar measure give $\int |u_z|^2>0$,
so $c_y>0$. In the convention of the doubling integral the
right-translation coefficient of $u_z$ at $b$ is therefore
$c_y\langle\rho(b)z,z\rangle$.

Choose a tensor $v$ of lower lowest-compact-type vectors for which
the equivariant coefficient map $J_v:C\to\mathcal S$ in the Fock
model is nonzero. Here is the nonvanishing for the particular model
used in this assertion. The normalization of its lowest-type map
sends one pair $(v,z)$ to the polynomial $\Delta_j$, a product of
powers of square minors in the plus and minus coordinates. Each
minor has a diagonal monomial with coefficient one, even when the
factor is evaluated in a different block. Their product is nonzero
in the polynomial ring. The finite vacuum is nonzero as well. In
the split half-model at the marked primes take in the two specified
frames the rectangular identity (the first $r$ rows form the
identity and the rest are zero). It lies in the required lattice,
satisfies the zero-block conditions and has all the tested unit
determinants equal to one. Projection to the marked finite
coordinates, with its conjugate on the dual half, gives this point
simultaneously. The remaining lattice vacua are nonzero. The
Schr\"odinger realization by Heisenberg translates of the initial
vacuum is injective on this half-model: disjoint residue cosets are
separated by the dual finite characters, and distinct characters
on a coset are linearly independent by orthogonality. Refining the
lattice deals with every locally constant compactly supported test.
Hence the chosen finite test survives. At the real places a nonzero
vacuum is killed by the annihilation operators. On its polynomials
in the creation operators the paired annihilators act by nonzero
multiples of partial differentiation. A derivative with the
multi-index of a nonzero term of highest total degree recovers
a nonzero multiple of the vacuum. Thus the vector attached to
$\Delta_j$ is nonzero. Fix such $v,z$, and put $\Phi=J_vz$.

The map $J_v$ intertwines $K$ and the Weil action on the upper
variable. Pair $J_vz$ and $J_vz'$ on the rational lattice quotient
of the theta model, using any covolume-one measure and fundamental
domain. The automorphy phases cancel; the Hermitian pairing is
$K$-invariant by unitarity of the Weil representation. By the
same Schur argument this pullback equals $c_v\langle z,z'\rangle$,
where $c_v>0$ because $J_vz$ is continuous and nonzero. The
Weil coefficient at $b$ equals $c_v\langle z,\rho(b)z\rangle$.
The archimedean factor of the doubling integral for the pair
$(\Phi,u_z)$ is consequently
\[
  c_v c_y \int_K |\langle z,\rho(b)z\rangle|^2\,db>0 .
\]
The integrand is continuous and its value at the identity is
$\|z\|^4>0$. Every nonempty open set has positive Haar measure
and $K$ is compact. The same $v,z$ work for all the covolume-one
choices and all Haar measures on $K$. In particular the scalar
component paired with this polynomial is nonzero for each $y\ne0$.
% END D87C2

% BEGIN D87C3
\paragraph{D87C3. Sphericity and the Euler product for this component.}
Keep the precise set $S$ of D87C1. The self-dual frame identities
at scale one identify both indicated lattices with the frame
lattices, by the two opposite inclusions. The projection in the
finite Weil step vacuum is supported on $Z$ and its conjugate.
A finite Heisenberg translation supported off $S$ has zero product
with both projections. It therefore acts on the new vacuum in
exactly the same way as on the original lattice vacuum. This
retains the lattice translation conditions, including doubled
translations at two. Real creation operators commute with these
finite translations, so the same holds for every polynomial
$\Phi$ of D87C2. At a principal-level element above $Z$ the
reductions are identity. The finite action on the step vacuum is
the inverse level character, which is one there; the lower lift
of the identity may itself be chosen to be the identity.
Commuting this action past the real polynomial gives principal
invariance of $\Phi$. The positive frame gives Cartan involution
$1$ on $V$, and the lower orthonormal frame gives the prescribed
Cartan involution on $W$. Hence $\Phi$ has all the original
finite and real properties in the occurrence hypothesis of the
doubling argument.

Suppose $y\ne0$ has spherical Hecke eigenvalues $s$ with the native
string property. At an unmarked finite place the level includes
the full hyperspecial group, with trivial character. Thus $u_z$
is fixed by that group. A spherical Hecke operator on $Y$ is
the sum of its finite right translations, with the same
normalization as on $y$. Applying $\langle z,-\rangle$ commutes
with that sum. Replacing coset representatives has no effect
by hyperspecial invariance. This proves that the nonzero scalar
form has the same eigenvalues $s$. At a split place the Satake
character determines its unordered parameters; at an inert
place the unramified base-change parameters are determined by
the relative Satake character. In particular the parameter
family $B$ witnessing the string for $s$ is a parameter family
for this very scalar form off $S$.

For $w$ over a place outside $S$, the uniformizer value of the
right-hand splitting character $\chi_V$ is the unit $z_w$ in
the string. All three factors in the member's splitting character
are unramified there: the fixed type character is unramified off
its designated primes, the twist is unramified off $Q(u_0)$ and
its conjugate, and $\eta$ is unramified. Define $c_w$ by subtracting
from $B_w$ the multiset $z_w q_w^{(t-1)/2-k}$, $0\le k<t$.
The string property ensures this subtraction is possible, gives
$|c_w|=r$ as a multiset cardinality, and bounds each remaining
root by $q_w^{(r-1)/2}$ in absolute value. Multiset subtraction
also works for repeated roots. At excluded places choose the
remainder arbitrarily. This gives the complete pole hypothesis
at exactly $S$, with both character arguments in the native
notation equal to $\chi_V$. Since $r+5\le t$, Y86N1 now gives
the simple pole at $s=t/2$ in the doubling variable for $u_z$.
% END D87C3

% BEGIN D87C4
\paragraph{D87C4. Applying definite nonvanishing to the original vector.}
For a member with linked realization, the geometric conditions
on the spaces and characters in the definite nonvanishing theorem
hold at every model index and tuple. The upper form is positive
definite at every real embedding. Its orthonormal bases supply
the real majorant; using their matrices at infinity and the identity
in the finite coordinates supplies the adelic choice. The lower
space has its given signature and positive rank $r$, and the upper
rank is $r+t$. The frame and the characteristic idempotent of $Z$
come from the level and step data. The host splits the places above
two. The character $\eta$ on the original Weil datum has the required
quadratic restriction: the sum of the two rank exponents is
$2r+t$, which is odd. Dividing the two splitting restrictions therefore gives the quadratic
restriction itself. The character agreement with the given Hecke
character gives continuity and rational triviality; the latter
character is unramified at every finite place and has exponent
$1/2$ at each specified infinity embedding. The same agreement
transports both facts to the character used in the Weil action. These checks use the same
rank, twist and $S$ as D87C3.

For any $y\ne0$ with the indicated eigenvalues and string, D87C1--3
supply the form, Fock vector and nonzero real integral in the
definite nonvanishing assertion $\mathrm{DownNe}$. Whenever that
assertion holds for these data there are a Schwartz test $f$ and a
lower argument $g$ for which the scalar lift of $u_z$ is nonzero.
Its value is the original integral $I_y(f,\langle z,-\rangle,g)$:
the product of the theta kernel with $u_z$ is invariant under the
upper rational group by the Weil automorphy law and D87C1. The
integral over the quotient therefore equals the integral over
the member's given fundamental domain with its recorded Haar
measure. This proves the implication from $\mathrm{DownNe}$ to
$\mathrm{NonVan}$ for each original tuple, with a coefficient
functional on the original vector even when the Hecke eigenspace
has dimension greater than one.
% END D87C4
